%# -*- coding: utf-8-unix -*-
\chapter*{追加文献}
\addcontentsline{toc}{chapter}{追加文献}
\markboth{追加文献}{追加文献}

% hanglist environment is defined globally in trans_number_baker.tex

我们下面给出了自本书首版印刷以来所出现的超越数\index{transcendental number}论新研究成果的精选列表; 该列表也包含了一些最初被遗漏的较早期的论文. 缩写与其他惯例与之前相同（参见第 130 页）.


\begin{hanglist}
    \item ANDERSON, M. Inhomogeneous linear forms in algebraic points of an elliptic function, \textit{Transcendence theory: advances and applications} (Academic Press, London, 1977), pp. 121-43.
    \item \hspace{1.5em} Linear forms in algebraic points of an elliptic function, Ph.D. dissertation (Nottingham, 1978).
    \item BAKER, A. Recent advances in transcendence theory, \textit{Trudy Mat. Inst. Steklov,} \textbf{132} (1973), 67-9.
    \item \hspace{1.5em} Some aspects of transcendence theory, \textit{Astirisque (Soc. Math. France),} \textbf{24-5} (1975), 169-75.
    \item \hspace{1.5em} The theory of linear forms in logarithms\index{linear forms in logarithms}, \textit{Transcendence theory: advances and applications} (Academic Press, London, 1977), pp. 1-27.
    \item \hspace{1.5em} Transcendence theory and its applications, \textit{J. Australian Math. Soc.} \textbf{25} (1978), 438-44.
    \item BAKER, A. and COATES, J. Fractional parts of powers of rationals, \textit{Math. P.C.P.S.,} \textbf{77} (1975), 269-79.
    \item BAKER, A. and STEWART, C. L. Further aspects of transcendence theory, \textit{Ast\'{e}risque (Soc. Math. France),} \textbf{41-2} (1977), 153-63.
    \item BAKER, R. C. On approximation with algebraic numbers\index{algebraic number} of bounded degree, \textit{Mathematika,} \textbf{23} (1976), 18-31.
    \item \hspace{1.5em} Spnndzuk's theorem and Hausdorff dimension\index{Hausdorff dimension}, \textit{ibid.} 184-97.
    \item \hspace{1.5em} Singular n-tuples and Hausdorff dimension\index{Hausdorff dimension}, \textit{Math. P.C.P.S.,} \textbf{81} (1977), 377-85.
    \item \hspace{1.5em} Dirichlet\index{Dirichlet, P. G. L.}'s theorem on Diophantine approximation, \textit{ibid.,} \textbf{83} (1978), 37-59.
    \item BERTRAND, D. Equations differentielles algebriques et nombres transcendants dans les domaines complexe et p-adique, Thesis (Paris VI, 1975); and related work \textit{C.R.} \textbf{279} (1974), 355-7; \textbf{282} (1976), 1399-401.
    \item \hspace{1.5em} Series d'Eisenstein et transcendance, \textit{Bull. Soc. Math. France,} \textbf{104} (1976), 309-21.
    \item \hspace{1.5em} Sous-groupes a un parametre p-adique de varietes de groupe, \textit{Invent. Math.} \textbf{40} (1977), 171-93.
    \item \hspace{1.5em} Sur les denominateurs des points rationnels des courbes elliptiques, \textit{Ast\'{e}risque (Soc. Math. France),} \textbf{41-2} (1977), 173-8.
    \item \hspace{1.5em} Algebraic values of p-adic elliptic functions\index{elliptic functions}, \textit{Transcendence theory: advances and applications} (Academic Press, London, 1977), pp. 149-58.
    \item \hspace{1.5em} A transcendence criterion for meromorphic functions\index{meromorphic function}, \textit{ibid.} pp. 187-93.
    \item \hspace{1.5em} Approximations Diophantiennesp-adiques sur les courbes elliptiques admettant une multiplication complexe, \textit{Compositio Math.} \textbf{37} (1978), 21-50.
    \item BERTRAND, D. and FLICKER, Y. Z. Linear forms on Abelian varieties over local fields, Acta Arith. 38 (1980), 47-61.
    \item Beukers\index{Beukers, F.}, F. On the generalized Ramanujan-Nagell equation I, II, Acta Arith. 38 (1981), 389-410; 39 (1981), 113-23.
    \item \hspace{1.5em} A note on the irrationality of  $\zeta(2)$  and  $\zeta(3)$ , Bull. London Math. Soc. 11 (1979), 268-72. Fractional parts of powers of rationals, Math. P.C.P.S. 90 (1981), 13-20.
    \item BIJLSMA, A. On the simultaneous approximation of a, b and ab, Compositio Math. 35 (1977), 99-111.
    \item Brownawell\index{Brownawell, W.}, W. D. Gelfond\index{Gelfond, A. O.}'s method for algebraic independence, \textit{Trans. Amer. Math. Soc.} 210 (1975), 1-26.
    \item \hspace{1.5em} A measure of linear independence for some exponential functions\index{exponential function}, \textit{Transcendence theory: advances and applications} (Academic Press, London, 1977), pp. 161-8.
    \item \hspace{1.5em} Some remarks on semi-resultants, ibid. pp. 205-10.
    \item \hspace{1.5em} On the Gelfond\index{Gelfond, A. O.}-Feldman\index{Feldman, N. I.} measure of algebraic independence, Compositio Math. 38 (1979), 355-68.
    \item BROWNAWELL, W. D. and KUBOTA, K. K. The algebraic independence of Weierstrass\index{Weierstrass, K.} functions and some related numbers\index{numbers}, \textit{Acta Arith.} 33 (1977), 111-49.
    \item Brownawell\index{Brownawell, W.}, W. D., and Waldschmidt\index{Waldschmidt, M.}, M. The algebraic independence of certain numbers\index{numbers} to algebraic powers, \textit{Acta Arith.} 32 (1977), 63-71.
    \item Bundschuh\index{Bundschuh, P.}, P. Irrationalität und Transzendenz gewisser Riehen, Math. Scand. 28 (1971), 226-32.
    \item \hspace{1.5em} Zu einem Transzendenzsatz der Herren Baron\index{Baron, G.} und Braune\index{Braune, E.}, J.M. 263 (1973), 183–8.
    \item \hspace{1.5em} Zum Franklin\index{Franklin, R.}-Schneiderschen Satz, J.M. 260 (1973), 108-18.
    \item \hspace{1.5em} Zur Approximation gewisser p-adischer algebraischer Zahlen durch rationale Zahlen, J.M. 265 (1974), 154-9.
    \item \hspace{1.5em} Zur simultanen Approximation von  $\beta_0, \ldots, \beta_{n-1}$  und  $\prod \alpha_p \beta_p$  durch algebraische Zahlen, J.M. 279 (1975), 99-117.
    \item \hspace{1.5em} Transzendenzmasse in Körpern formaler Laurentreihen, J.M. 300 (1978), 411-32.
    \item Bundschuh\index{Bundschuh, P.}, P. and Wallisser\index{Wallisser, R.}, R. Algebraische Unabhängigkeit, P-adischer Zahlen, M.A. 221 (1976), 243-9.
    \item Chudnovsky\index{Chudnovsky, G. V.}, G. V. A measure of mutual transcendence for certain number classes (in Russian), D.A.N. 18 (1974), 771-4.
    \item \hspace{1.5em} Baker\index{Baker, R. C.}\index{Baker, A.}'s method in the theory of transcendental numbers\index{transcendental number}\index{numbers} (in Russian), U.M. 31 (1976), 281-2.
    \item \hspace{1.5em} The Gelfond\index{Gelfond, A. O.}-Baker\index{Baker, R. C.}\index{Baker, A.} method in problems of Diophantine approximation, Colloq. Math. Soc. János Bolyai, vol. 13 (1976), pp. 19-30.
    \item \hspace{1.5em} A new method for the investigation of arithmetical properties of analytic functions, Ann. Math. 109 (1979), 353-76.
    \item Cijsouw\index{Cijsouw, P. L.}, P. L. Transcendence measures of exponentials and logarithms of algebraic numbers\index{logarithms of algebraic numbers}\index{algebraic number}, \textit{Compositio Math.} 28 (1974), 163-78.
    \item \hspace{1.5em} Transcendence measures of certain numbers whose transcendency was proved by A. Baker\index{Baker, R. C.}\index{Baker, A.}, \textit{ibid}, 179-94.
    \item \hspace{1.5em} On the approximability of the logarithms of algebraic numbers\index{logarithms of algebraic numbers}\index{algebraic number}, Sém. Delange-Pisot-Poitou 16 (1975), No. 19.
    \item \hspace{1.5em} On the simultaneous approximation of certain numbers\index{numbers}, Duke Math. J. 42 (1975), 249-57.
    \item \hspace{1.5em} A transcendence measure for  $\pi$ , Transcendence theory: advances and applications (Academic Press, London, 1977), pp. 93-100.
    \item CIJSOUW, P. L. and TIJDEMAN, R. An auxiliary result in the theory of transcendental numbers\index{transcendental number}\index{numbers} II, \textit{Duke Math. J.} 42 (1975), 239-47.
    \item CIJSOUW, P. L. and WALDSCHMIDT, M. Linear forms and simultaneous approximations, Compositio Math. 34 (1977), 173-97.
    \item COATES, J. and LANG, S. Diophantine approximation on Abelian varieties with complex multiplications\index{complex multiplication}, \textit{Invent. Math.} 34 (1976), 129-33.
    \item Cusick\index{Cusick, T. W.}, T, W. Effective lower bounds for some linear forms, Trans. Amer. Math. Soc. 222 (1976), 289-301.
    \item Davis\index{Davis, M.}\index{Davis, C. S.}, C. S. Rational approximations to e, J. Australian Math. Soc. 25 (1978), 497-502.
    \item DAVISON, J. L. A series and its associated continued fraction, Proc. Amer. Math. Soc. 63 (1977), 29-32; and related work with W. W. Adams\index{Adams, W. W.}, ibid. 65 (1977), 194-8.
    \item DEMJANENKO, V. A. Tate height and the representation of numbers\index{numbers} by binary forms (in Russian), I.A.N. 38 (1974), 459-70.
    \item DUBOIS, E. Application de la méthode de W. M. Schmidt\index{Schmidt, W. M.} à l'approximation de nombres algébriques dans un corps de fonctions de caractéristique zéro, C.R. 284 (1977), 1527-30.
    \item Dubois\index{Dubois, E.}, E. and Rhin\index{Rhin, G.}, G. Approximations rationnelles simultanées de nombres algébriques réels et de nombres algébriques p-adiques, Astérisque (Soc. Math. France), 24-5 (1975), 211-27.
    \item \hspace{1.5em} Sur la majoration de formes linéaires à coefficients algébriques réels et p-adiques. Demonstration d'une conjecture de K. Mahler\index{Mahler, K.}, C.R. 282 (1976), 1211-14.
    \item DURAND, A. Quatre problèmes de Mahler\index{Mahler, K.} sur la fonction ordre d'un nombre transcendant, \textit{Bull. Soc. Math. France}, \textbf{102} (1974), 365-77; and related work \textit{C.R.} \textbf{280} (1975), 309-11, 1085-8.
    \item \hspace{1.5em} Note on rational approximations of the exponential function\index{exponential function} at rational points, Bull. Australian Math. Soc. 14 (1976), 449-55.
    \item \hspace{1.5em} Indépendance algébrique de nombres complexes et critère de transcendance, Compositio Math. 35 (1977), 259-67.
    \item ELIANU, JEAN. Sur le nombre  $[a, a + \beta, a + 2\beta, ..., a + n\beta, ...]$ . Rev. Roumaine Math. Pures Appl. 20 (1975), 1061-71.
    \item Endos, P. and Shorey\index{Shorey, T. N.}, T. N. On the greatest prime factor of  $2^p-1$  for a prime p, and other expressions, \textit{Acta Arith.} 30 (1976), 257-65.
    \item FELDMAN, N. I. The periods of elliptic functions\index{elliptic functions} (in Russian), Acta Arith. 24 (1974), 477-89.
    \item \hspace{1.5em} Estimates of linear forms in the logarithms of algebraic numbers\index{logarithms of algebraic numbers}\index{algebraic number}, and some applications (in Russian), Current problems in analytic number theory (Minsk, 1974), pp. 244-68.
    \item \hspace{1.5em} Rational approximations of algebraic numbers\index{algebraic number}, Colloq. Math. Soc. János Bolyai, vol. 13 (1976), pp. 31-9.
    \item \hspace{1.5em} Approximation of number  $\pi$  by algebraic numbers\index{algebraic number} from special fields, J. Number Theory, 9 (1977), 48-60.
    \item FLICKER, Y. Z. On p-adic G-functions\index{G-functions}, J. London Math. Soc. 15 (1977), 395-402.
    \item \hspace{1.5em} Transcendence theory over local fields, Ph.D. dissertation (Cambridge, 1978).
    \item \hspace{1.5em} Algebraic independence by a method of Mahler\index{Mahler, K.}, J. Australian Math. Soc. 27 (1979), 173-88.
    \item FRANKLIN, R. The transcendence of linear forms in  $\omega_1$ ,  $\omega_2$ ,  $\eta_1$ ,  $\eta_2$ ,  $2\pi i$ ,  $\log \gamma$ , Acta Arith. 26 (1975), 197–206; but see also 31 (1976), 143–52.
    \item GALOČKIN, A. I. Lower bounds of polynomials in the values of a certain class of analytic functions (in Russian), \textit{Mat. Sb.} 95 (1974), 396-417.
    \item \hspace{1.5em} The approximation of the parameters of certain hypergeometric functions\index{hypergeometric function} (in Russian), \textit{Mat. Zametki}, 17 (1975), 103-12; and related work, \textit{ibid}. 18 (1975), 541-52; 20 (1976), 35-45; \textit{V.M.}, no. 6 (1978), 25-32.
    \item Gerritzen\index{Gerritzen, L.}, L. Ein p-adischer Beweis für die Irrationalität der Nullstellen von Besselfunktionen, M.A. 226 (1977), 253-5.
    \item GROSSMAN, E. H. Units and discriminants of algebraic number\index{algebraic number} fields, Comm. Pure Appl. Math. 27 (1974), 741-7.
    \item Györy, K. Sur l'irreductibilité d'une classe des polynômes I, II, Publ. Math. Debrecen, 18 (1971), 289-307; 19 (1972), 293-326.
    \item \hspace{1.5em} On polynomials with integer coefficients and given discriminant I, II, III, IV, V, Acta Arith. 23 (1973), 419-26; Publ. Math. Debrecen 21 (1974), 125-44; 23 (1976), 141-65; 25 (1978), 155-67; Acta Math. Hungar. 32 (1978), 175-90.
    \item \hspace{1.5em} Polynomials with given discriminant, Colloq. Math. Soc. János Bolyai, vol. 13 (1974), 65-78.
    \item \hspace{1.5em} Sur une classe des corps de nombres algébriques et ses applications, \textit{Publ.}Math. Debrecen 22 (1975), 151-75.
    \item \hspace{1.5em} Représentation des nombres entiers par des formes binaires, \textit{ibid}. 24 (1977), 363-75.
    \item Györy, K. and Lovász, L. Representation of integers by norm forms II, Publ. Math. Debrecen, 17 (1970), 173-81; Part I, by K. Györy, ibid. 16 (1969), 253-63.
    \item Györy, K. and Papp\index{Papp, Z. Z.}, Z. Z. Effective estimates for the integer solutions of norm form and discriminant form equations, \textit{Publ. Math. Debrecen} 25 (1978), 311-25; and similar titles to appear.
    \item Györy, K., Tijdeman\index{Tijdeman, R.}, R. and Voorhoeve\index{Voorhoeve, M.}, M. On the equation  $1^k + 2^k + \dots + x^k = y^z$ , Acta Arith. (to appear).
    \item HARASE, T. The transcendence of  $e\index{transcendence of e}^{\alpha\omega+\beta(2\pi i)}$ , J. Fac. Sci. Univ. Tokyo 21 (1974), 279-85.
    \item \hspace{1.5em} On the linear form of transcendental numbers\index{transcendental number}\index{numbers}  $a_1\omega + a_2\pi + a_3 \log a_4$ , \textit{ibid}. 23 (1976), 435-52.
    \item KLEIMAN, H. On the Diophantine equation f(x, y) = 0, J.M. 287 (1976), 124-31.
    \item Korov, S. V. The law of the iterated logarithm for binary forms with algebraic coefficients (in Russian), Dokl. Akad. Nauk BSSR 17 (1973), 591-4.
    \item \hspace{1.5em} The greatest prime factor of a polynomial (in Russian), Mat. Zametki, 13 (1973), 515-22; = 13 (1973), 313-17.
    \item \hspace{1.5em} The Thue\index{Thue, A.}-Mahler\index{Mahler, K.} equation in relative fields (in Russian), Acta Arith. 27 (1975), 293-315.
    \item \hspace{1.5em} Über die maximale Norm der Idealteiler des Polynoms  $ax^m + \beta y^n$  mit den algebraischen Koeffizienten, \textit{ibid.} 31 (1976), 219-30.
    \item Kubota\index{Kubota, K. K.}, K. K. On the algebraic independence of holomorphic solutions of certain functional equations and their values, M.A. 227 (1977), 9-50.
    \item \hspace{1.5em} On a transcendence problem of K. Mahler\index{Mahler, K.}, Canadian J. Math. 29 (1977), 638-47.
    \item \hspace{1.5em} Linear functional equations and algebraic independence, Transcendence theory: advances and applications (Academic Press, London, 1977), pp. 227-9.
    \item LANG, S. Diophantine approximation on toruses, Amer. J. Math. 86 (1964), 521-33.
    \item Division points of elliptic curves\index{elliptic curve} and abelian functions\index{Abelian function} over number fields, \textit{ibid}. 97 (1972), 124-32.
    \item Higher dimensional Diophantine problems, Bull. Amer. Math. Soc. 80 (1974), 779-87.
    \item Diophantine approximation on abelian varieties with complex multiplication\index{complex multiplication}, Advances in Math. 17 (1975), 281-336.
    \item Elliptic curves\index{elliptic curve}. Diophantine analysis\index{Diophantine analysis} (Springer-Verlag, Berlin, Heidelberg, 1978).
    \item LEVEQUE, W. J. On the equation  $y^m = f(x)$ , Acta Arith. 9 (1964), 209-19.
    \item LOXTON, J. H. and VAN DER POORTEN, A. J. On algebraic functions satisfying a class of functional equations, \textit{Aequationes Math.} 14 (1976), 413-20; and related work \textit{ibid.} 15 (1977), 114-15.
    \item \hspace{1.5em} On the growth of recurrence sequences, Math. P.C.P.S. 81 (1977), 369-76.
    \item \hspace{1.5em} Arithmetic properties of certain functions in several variables I, II, III, J. Number Theory, 9 (1977), 87-106; J. Australian Math. Soc. (to appear); Bull. Australian Math. Soc. 16 (1977), 15-47.
    \item \hspace{1.5em} Transcendence and algebraic independence by a method of Mahler\index{Mahler, K.}, Transcendence theory: advances and applications (Academic Press, London, 1977), pp. 211-26.
    \item Mahler\index{Mahler, K.}, K. Arithmetische Eigenschaften der Lösungen einer Klasse von Funktionalgleichungen, M.A. 101 (1929), 342-66.
    \item \hspace{1.5em} Über das Verschwinden von Potenzreihen mehrerer Veränderlichen in speziellen Punktfolgen, M.A. 103 (1930), 573-87.
    \item \hspace{1.5em} Arithmetische Eigenschaften einer Klasse transzendental-transzendenter Funktionen, M.Z. 32 (1930), 545–85.
    \item \hspace{1.5em} Lectures on transcendental numbers\index{transcendental number}\index{numbers}, Proc. Symposia Pure Math., vol. 20 (Amer. Math. Soc., 1971), pp. 248-74.
    \item \hspace{1.5em} The classification of transcendental numbers\index{transcendental number}\index{numbers}, Proc. Symposia Pure Math., vol. 24 (Amer. Math. Soc., 1973), pp. 175-9.
    \item \hspace{1.5em} A p-adic analogue to a theorem of J. Popken\index{Popken, J.}, J. Australian Math. Soc. 16 (1973), 176-84.
    \item \hspace{1.5em} On rational approximations of the exponential function\index{exponential function} at rational points, Bull. Australian Math. Soc. 10 (1974), 325-35.
    \item \hspace{1.5em} On the transcendency of the solutions of a special class of functional equations, \textit{ibid.} 13 (1975), 389-410; corrigendum 14 (1976), 477-8.
    \item \hspace{1.5em} A necessary and sufficient condition for transcendency, Math. Comput. 29 (1975), 145-53.
    \item \hspace{1.5em} On a paper by A. Baker\index{Baker, R. C.}\index{Baker, A.} on the approximation of rational powers of e, Acta Arith. 27 (1975), 61-87.
    \item \hspace{1.5em} On a class of transcendental decimal fractions, Comm. Pure Appl. Math. 29 (1976), 717-25; and a similar title in Number theory and algebra (Academic Press, New York, 1977), pp. 209-14.
    \item \hspace{1.5em} Lectures on transcendental numbers\index{transcendental number}\index{numbers}, Lecture Notes in Mathematics, vol. 546 (Springer-Verlag, Berlin, 1976).
    \item MAKAROV, Yu. N. On the estimates of the measure of linear independence for the values of E-functions\index{E-functions} (in Russian), V.M., no. 2 (1978), 3-12.
    \item Mandelbrojt\index{Mandelbrojt, S.}, S. Sur la décroissance des coefficients supposés rationnels d'une fonction entière ne s'annulant pas en e, C.R. 278 (1974), 921-4.
    \item MASSER, D. W. Elliptic functions\index{elliptic functions} and transcendence, Lecture Notes in Mathematics, vol. 437 (Springer-Verlag, Berlin, 1975).
    \item Linear forms in algebraic points of Abelian functions\index{Abelian function} I, II, III, Math. P.C.P.S. 77 (1975), 499-513; 79 (1976), 55-70; Proc. London Math. Soc. 33 (1976), 549-64.
    \item On the periods of Abelian functions\index{Abelian function} in two variables, \textit{Mathematika}, 22 (1975), 97-107.
    \item A note on a paper of Franklin\index{Franklin, R.}, Acta Arith, 31 (1976), 143-52.
    \item Division fields of elliptic functions\index{elliptic functions}, Bull. London Math. Soc. 9 (1977), 49-53.
    \item The transcendence of certain quasi-periods associated with Abelian functions\index{Abelian function} in two variables, \textit{Compositio Math.} 35 (1977), 239-58.
    \item The transcendence of definite integrals of algebraic functions, Astérisque (Soc. Math. France), 41-2 (1977), 231-8.
    \item Some vector spaces associated with two elliptic functions\index{elliptic functions}, \textit{Transcendence theory: advances and applications} (Academic Press, London, 1977), pp. 101-19.
    \item A note on Abelian functions\index{Abelian function}, ibid., pp. 145-7.
    \item Diophantine approximation and lattices with complex multiplication\index{complex multiplication}, \textit{Invent. Math.} 45 (1978), 61-82.
    \item MEYER, Y. Nombres transcendants et répartition modulo 1, Bull. Soc. Math. France, 25 (1971), 143-9.
    \item MIGNOTTE, M. Une généralisation d'un théorème de Cugiani\index{Cugiani, M.}-Mahler\index{Mahler, K.}, Acta Arith. 22 (1972), 57-67.
    \item \hspace{1.5em} Demonstration probabiliste d'un lemme combinatoire pour l'approximation diophantienne des nombres algébriques, Collog. Math. 30 (1974), 109-19.
    \item \hspace{1.5em} Approximations rationelles de  $\pi$  et quelques autres nombres, \textit{Bull. Soc. Math. France}, 37 (1974), 121-32.
    \item \hspace{1.5em} Quelques problèmes d'effectivite en theorie des nombres, Thesis (Paris XIII, 1974).
    \item \hspace{1.5em} Quelques remarques sur l'approximation rationnelle des nombres algébriques, J.M. 269 (1974), 341-7.
    \item \hspace{1.5em} Approximation des nombres algébriques par certaines suites de rationnels, Rend. Circ. Mat. Palermo, 24 (1975), 87-124.
    \item \hspace{1.5em} A note on linear recursive sequences, J. Australian Math. Soc. 20 (1975), 242-4
    \item MIGNOTTE, M. and WALDSCHMIDT, M. Approximation des valeurs de fonctions transcendantes, \textit{Indag. Math.} 37 (1975), 213-23.
    \item \hspace{1.5em} Approximation simultanée de valeurs de la fonction exponentielle, Compositio Math. 34 (1977), 127-39.
    \item \hspace{1.5em} Linear forms in two logarithms and Schneider\index{Schneider, Th.}'s method, M.A. 231 (1978), 241-67.
    \item Moreno\index{Moreno, C. J.}, C. J. The values of exponential polynomials\index{exponential polynomials} at algebraic points I, II, Trans. Amer. Math. Soc. 186 (1973), 17-31; Diophantine approximation and its applications (Academic Press, New York, 1973), pp. 111-28.
    \item NESTERENKO, Ju. V. Estimates for the orders of the zeros of analytic functions of a certain class and their applications in the theory of transcendental numbers\index{transcendental number}\index{numbers} (in Russian), D.A.N. 205 (1972), 292-5; and a similar title I.A.N. 41 (1977), 253-84.
    \item \hspace{1.5em} An order function for almost all numbers\index{numbers} (in Russian), Mat. Zametki, 15 (1974), 405-14.
    \item NURMAGOMEDOV, M. S. and CHIRSKY, V. G. On arithmetical properties of values of hypergeometric functions\index{hypergeometric function} (in Russian), V.M. 28, no. 2 (1973), 38-45; see also [V.G. Chirsky\index{Chirsky, V. G.}] no. 5 (1978), 3-8.
    \item VAN DER POORTEN, A. J. Hermite\index{Hermite, Ch.} interpolation and p-adic exponential polynomials\index{exponential polynomials}, J. Australian Math. Soc. 21 (1976), 12-26.
    \item \hspace{1.5em} On Baker\index{Baker, R. C.}\index{Baker, A.}'s inequality for linear forms in logarithms\index{linear forms in logarithms}, Math. P.C.P.S. 80 (1976), 233-48.
    \item \hspace{1.5em} Effectively computable bounds for the solutions of certain Diophantine equations\index{Diophantine equations}, Acta Arith. 33 (1977), 195-207.
    \item \hspace{1.5em} Linear forms in logarithms\index{linear forms in logarithms} in the p-adic case, Transcendence theory: advances and applications (Academic Press, London, 1977), pp. 29-57.
    \item VAN DER POORTEN, A. J. and LOXTON, J. H. Computing the effectively computable bound in Baker\index{Baker, R. C.}\index{Baker, A.}'s inequality for linear forms in logarithms\index{linear forms in logarithms}, Bull. Australian Math. Soc. 15 (1976), 33-57; and related work, 16 (1977), 83-98; corrigendum, 17 (1977), 151-5.
    \item RAMACHANDRA, K., SHOREY, T. N. and TIJDEMAN, R. On Grimm's problem relating to factorisation of a block of consecutive integers I, II, J.M. 273 (1975), 109-24; 288 (1976), 192-201.
    \item REYSSAT, E. Mesures de transcendance de nombres lies aux fonctions elliptiques, Thesis (Paris VI, 1977).
    \item Salihov\index{Salihov, V. H.}, V. H. The algebraic independence of the values of \textit{E}-functions that satisfy first order linear differential equations (in Russian), \textit{Mat. Zametki}, 13 (1973), 29-40; and related work, \textit{D.A.N.} 235 (1977), 30-3.
    \item SCHINZEL, A. and TIJDEMAN, R. On the equation  $y^m = f(x)$ , Acta Arith. 31 (1976), 199-204.
    \item Schlickewei\index{Schlickewei, H. P.}, H. P. On the fractional parts of the sum of powers of rational numbers\index{numbers}, \textit{Mathematika}, 22 (1975), 154-5.
    \item \hspace{1.5em} Die p-adische Verallgemeinerung des Satzes von Thue\index{Thue, A.}-Siegel\index{Siegel, C. L.}-Roth\index{Roth, K. F.}-Schmidt\index{Schmidt, W. M.}, J.M. 288 (1976), 86-105.
    \item \hspace{1.5em} Linearformen mit algebraischen koeffizienten, \textit{Manuscripta Math.} 18 (1976), 147–85; and related work, \textit{Acta Arith.} 31 (1976), 389–98; 33 (1977), 183–5.
    \item \hspace{1.5em} On norm form equations\index{norm form equation}, J. Number Theory, 9 (1977), 370-80; and related work, ibid. 381-92; Astérisque (Soc. Math. France), 41-2 (1977), 267-71.
    \item SCHMIDT, W. M. Approximation to algebraic numbers\index{algebraic number}, Enseignement Math. 17 (1971), 187-253.
    \item \hspace{1.5em} Zur Methode von Stepanov\index{Stepanov, S. A.}, Acta Arith. 24 (1973), 347-67; and related work, J. Number Theory, 6 (1974), 448-80.
    \item \hspace{1.5em} Simultaneous approximation to algebraic numbers\index{algebraic number} by elements of a number field, \textit{Monatsh. Math.} 79 (1975), 55-66.
    \item \hspace{1.5em} Rational approximation to solutions of linear differential equations with algebraic coefficients, \textit{Proc. Amer. Math. Soc.} 53 (1975), 285-9.
    \item \hspace{1.5em} Equations over finite fields. An elementary approach, Lecture Notes in Mathematics, vol. 536 (Springer-Verlag, Berlin, 1976).
    \item \hspace{1.5em} On Osgood\index{Osgood, C. F.}'s effective Thue\index{Thue, A.} theorem for algebraic functions. Comm. Pure Appl. Math. 29 (1976), 749-63; and related work Acta Arith. 32 (1977), 275-96; J. Australian Math. Soc. 25 (1978), 385-422.
    \item Schneider\index{Schneider, Th.}, Th. Eine Bemerkung zu einem Satz von C. L. Siegel\index{Siegel, C. L.}, Comm. Pure Appl. Math. 29 (1976), 775-82.
    \item Shidlovsky\index{Shidlovsky, A. B.}, A. B. The arithmetic properties of the values of analytic functions (in Russian). \textit{Trudy Mat. Inst. Steklov}, \textbf{132} (1973), 169–202.
    \item SHOREY, T. N. On the sum  $\sum_{k=1}^{3} |2^{\pi^k} a_k|$ ,  $a_k$  algebraic numbers\index{algebraic number}, J. Number Theory, 6 (1974), 248-60.
    \item On linear forms in the logarithms of algebraic numbers\index{logarithms of algebraic numbers}\index{algebraic number}, Acta Arith. 30 (1976), 27-42.
    \item SHOREY, T. N. and TIJDEMAN, R. New applications of Diophantine approximations to Diophantine equations\index{Diophantine equations}, \textit{Math. Scand.} 39 (1976), 5-18.
    \item \hspace{1.5em} On the greatest prime factors of polynomials at integer points, Compositio Math. 33 (1976), 187-95.
    \item SHOREY, T. N., VAN DER POORTEN, A. J., TIJDEMAN, R. and SCHINZEL, A. Applications of the Gelfond\index{Gelfond, A. O.}-Baker\index{Baker, R. C.}\index{Baker, A.} method to Diophantine equations\index{Diophantine equations}, \textit{Transcendence theory: advances and applications} (Academic Press, London, 1977), pp. 59-77.
    \item ŠMELEV\index{Smelev, A. A.}, A. A. On a method of A. O. Gelfond\index{Gelfond, A. O.} in the theory of transcendental numbers\index{transcendental number}\index{numbers} (in Russian), Mat. Zametki, 10 (1971), 415-26; see also 11 (1972), 635-44.
    \item \hspace{1.5em} Simultaneous approximations of the values of the exponential function\index{exponential function} at nonalgebraic points (in Russian), V.M. 27, no. 4 (1972), 25-33; see also no. 6 (1972), 5-14, and Ukrain. Mat. Z. 27 (1975), 555-63.
    \item \hspace{1.5em} A criterion for the algebraic dependence of transcendental numbers\index{transcendental number}\index{numbers} (in Russian), Mat. Zametki, 16 (1974), 553-62.
    \item SPRINDŽUK\index{Sprindzhuk, V. G.}, V. G. Square-free divisors of polynomials and class numbers of algebraic number\index{algebraic number} fields (in Russian), Acta Arith. 24 (1973), 143-9.
    \item \hspace{1.5em} Algebraic number\index{algebraic number} fields with large class number (in Russian), I.A.N. 38 (1974), 971-82; = 8 (1974), 967-78.
    \item \hspace{1.5em} The metric theory of Diophantine approximations (in Russian), Current problems of analytic number theory (Minsk, 1974), pp. 178-88; see also Trudy Mat. Inst. Steklov, 132 (1973), 137-42.
    \item \hspace{1.5em} An effective analysis of the Thue\index{Thue, A.} and Thue\index{Thue, A.}-Mahler\index{Mahler, K.} equations (in Russian), Current problems of analytic number theory (Minsk, 1974), pp. 199-222.
    \item \hspace{1.5em} Representation of numbers\index{numbers} by the norm forms with two dominating variables, J. Number Theory, 6 (1974), 481-6.
    \item \hspace{1.5em} A hyperelliptic Diophantine equation and class numbers\index{numbers} (in Russian), Acta Arith. 30 (1976), 95-108.
    \item SPRINDŽUK\index{Sprindzhuk, V. G.}, V. G. and KOTOV, S. V. An effective analysis of the Thue\index{Thue, A.}-Mahler\index{Mahler, K.} equation in relative fields (in Russian), Dokl. Akad. Nauk. BSSR, 17 (1973), 393-5.
    \item \hspace{1.5em} The approximation of algebraic numbers\index{algebraic number} by algebraic numbers\index{algebraic number} of a given field (in Russian), \textit{ibid}. 20 (1976), 581-4.
    \item SRINIVASAN, S. On algebraic approximations to  $2^{n^k}(k=1,2,3,...)$ , India J. Pure Appl. Math. 5 (1974), 513-23.
    \item STARE, H. M. Some effective cases of the Brauer\index{Brauer, R.}\index{Brauer, A.}-Siegel\index{Siegel, C. L.} theorem, Invent. Math. 23 (1974), 135-52; see also Bull. Amer. Math. Soc. 81 (1975), 961-72
    \item \hspace{1.5em} On complex quadratic fields with class-number two, Math. Comp. 29 (1975) 289-302.
    \item STEWART, C. L. The greatest prime factor of  $a^n b^n$ , Acta Arith. 26 (1975), 427-33.
    \item \hspace{1.5em} Divisor properties of arithmetical sequences, Ph.D. dissertation (Cambridge, 1976).
    \item \hspace{1.5em} On divisors of Fermat, Fibonacci, Lucas and Lehmer numbers\index{Lucas and Lehmer numbers}, Proc. London Math. Soc. 35 (1977), 425-47.
    \item \hspace{1.5em} Primitive divisors of Lucas and Lehmer numbers\index{Lucas and Lehmer numbers}, Transcendence theory: advances and applications (Academic Press, 1977), pp. 79-92.
    \item \hspace{1.5em} A note on the Fermat equation\index{Fermat equation}, Mathematika, 24 (1977), 130-2.
    \item Algebraic integers\index{algebraic integer} whose conjugates lie near the unit circle, Bull. Soc. Math. France, 106 (1978), 169-76.
    \item STOLARSKY, K. B. Algebraic numbers\index{algebraic number} and Diophantine approximation (Marcel Dekker, New York, 1974).
    \item TIJDEMAN, R. Old and new in number theory, Nieuw Arch. Wisk. 20 (1972), 20-30.
    \item \hspace{1.5em} On the maximal distance between integers composed of small primes, Compositio Math. 28 (1974), 159-62; for applications see joint work with H. G. Meijer, ibid. 29 (1974), 265-86.
    \item \hspace{1.5em} Applications of the Gelfond\index{Gelfond, A. O.}-Baker\index{Baker, R. C.}\index{Baker, A.} method to rational number theory, Collog. Math. Soc. János Bolyai, vol. 13 (1974), pp. 399-416.
    \item \hspace{1.5em} Some applications of Baker\index{Baker, R. C.}\index{Baker, A.}'s sharpened bounds to Diophantine equations\index{Diophantine equations}, Sém. Delange-Pisot-Poitou, 16 (1975), No. 24.
    \item \hspace{1.5em} Hilbert\index{Hilbert, D.}'s seventh problem\index{Hilbert's seventh problem}: On the Gelfond\index{Gelfond, A. O.}-Baker\index{Baker, R. C.}\index{Baker, A.} method and its applications, Proc. Symposia Pure Math., vol. 28 (Amer. Math. Soc., 1976), pp. 241-68.
    \item \hspace{1.5em} On the equation of Catalan, Acta Arith. 29 (1976), 197-209.
    \item VÄÄNÄNEN, K. On a conjecture of Mahler\index{Mahler, K.} concerning the algebraic independence of the values of some E-functions\index{E-functions}, Ann. Acad. Sci. Fenn. Ser. AI, no. 512 (1972); and similar titles, ibid. no. 536 (1973), and Math. 1 (1975), pp. 93-109, 183-94.
    \item \hspace{1.5em} On lower estimates for linear forms involving certain transcendental numbers\index{transcendental number}\index{numbers}, Bull. Australian Math. Soc. 14 (1976), 161-79.
    \item \hspace{1.5em} On lower bounds for polynomials in the values of E-functions\index{E-functions}, Manuscripta Math. 21 (1977), 173-80; see also J.M. 296 (1977), 205-11.
    \item VOORHOEVE, M., GYÖRY, K. and TIJDEMAN, R. On the diophantine equation  $1^k + 2^k + \ldots + x^k + R(x) = y^x Acta Math.$  143 (1979), 1-8.
    \item WALDSCHMIDT, M. Propriétés arithmétiques de fonctions méromorphes, C.R. 273 (1971), 554-7.
    \item \hspace{1.5em} Dimension algébrique de sous-groupes analytiques de variétés abéliennes, C.R. 274 (1972), 1681-3.
    \item \hspace{1.5em} Utilisation de la méthode de Baker\index{Baker, R. C.}\index{Baker, A.} dans des problèmes d'indépendance algébrique, C.R. 275 (1972), 1215-17.
    \item \hspace{1.5em} Approximation par des nombres algébriques des zeros de séries entières à coefficients algébriques, C.R. 279 (1974), 793-6.
    \item \hspace{1.5em} Initiation aux nombres transcendants, Enseignement Math. 20 (1974), 53-85. A propos de la méthode de Baker\index{Baker, R. C.}\index{Baker, A.}, Bull. Soc. Math. France, 37 (1974), 181-92.
    \item \hspace{1.5em} Indépendance algébrique par la méthode de G. V. Chudnovsky\index{Chudnovsky, G. V.}, Sém. Delange-Pisot-Poitou, 16 (1975), no. 8.
    \item \hspace{1.5em} Images de points algébriques par un sous-groupe analytique d'une variété de groupe, C.R. 281 (1975), 855-8.
    \item \hspace{1.5em} Dimension algébrique de sous-groupes analytiques de variétés de groupe, Ann. Inst. Fourier, Grenoble, 25 (1975), 23-33.
    \item \hspace{1.5em} Some topics in transcendental number\index{transcendental number} theory, Colloq. Math. Soc. Janós Bolyai, vol. 13 (1976), pp. 417-27.
    \item \hspace{1.5em} Rapport sur la transcendance, Astérisque (Soc. Math. France), 41-2 (1977), 127-34.
    \item \hspace{1.5em} On functions of several variables having algebraic Taylor coefficients, Transcendence theory: advances and applications (Academic Press, London, 1977), pp. 169-86; see also articles in Lecture Notes in Mathematics, vol. 524 (1976), vol. 578 (1977) (Springer-Verlag, Berlin).
    \item Lea travaux de G. V. Chudnovsky\index{Chudnovsky, G. V.} sur les nombres trp.nscendants, \textit{ibid.} vol. 567 (1977), pp. 274-92.
    \item Transcendence measures for exponentials and logarithms, \textit{J. Australian Math.Soc.} 25 (1978), 445-65.
    \item Simultaneous approximations of numbers connected with the exponential function\index{exponential function}, \textit{ibid.} 466-78.
    \item A lower bound for linear forms in logarithms\index{linear forms in logarithms}, \textit{Ada Arilh.} 37 (1980), 257-83.
    \item WALLISSER, R. Uber die arithmetische Natur der Werte der Losungen einer Funktionalgleichung von H. Poincare, \textit{Acta Arith.} 25 (1974), 81-92.
    \item WUSTHOLZ\index{Wüstholz, G.}, G. Zum Franklin\index{Franklin, R.}-Schneiderschen Satz, \textit{Manuscripta Math.} 20 (1977), 335-54.
    \item \hspace{1.5em} Linearformen in Logarithmen von [/-Zahlen mit ganzzahligen Koefnzienten, \textit{J.M.} 300 (1978), 138-50.
\end{hanglist}
